% !TeX root = ../main.tex
\subsection{تمرین \secnum{۵.۱۰}، اندازه‌گذاری \lr{pMOS Header}}
\label{sec:assignment-power-header}

\textbf{خواسته.}
در یک مدار \lr{Power Gating} با فرایند \qty{65}{\nano\metre}،
عرض کلید \lr{pMOS Header} را تعیین کنید تا افزایش تأخیر کمتر از
۲ درصد باشد. حاصل‌ضرب تقریبی مقاومت روشن کلید در عرض آن
\(r_w=\qty{2.5}{\kilo\ohm\micro\metre}\) و جریان فعال بلوک
\(I_{\mathrm{on}}=\qty{100}{\milli\ampere}\) است
\cite[تمرین \secnum{۵.۱۰}، صفحه ۲۱۰]{weste}.

\textbf{حل.}
با فرض تغذیه \(V_{DD}=\qty{1}{\volt}\) از مثال ۵.۵ کتاب،
از همان تقریب استفاده می‌کنیم که افزایش نسبی کوچک تأخیر را
با افت نسبی تغذیه برابر می‌گیرد. بنابراین در این مدل، برای
افزایش کمتر از ۲ درصد باید \(\Delta V<0.02V_{DD}=\qty{20}{\milli\volt}\)
باشد \cite[مثال ۵.۵، صفحات ۱۹۸ و ۱۹۹ و رابطه \secnum{۴.۲۹}،
صفحه ۱۶۲]{weste}.

با \(W\) به‌عنوان عرض مؤثر کلید و \(R_h\) به‌عنوان مقاومت روشن،
داریم \(R_h\approx r_w/W\) و \(\Delta V\approx I_{\mathrm{on}}R_h\).
در مدل مقاومتی، شرطِ افت به کران‌های زیر می‌انجامد.
\[
  \begin{aligned}
    R_h &< \frac{\qty{20}{\milli\volt}}{\qty{100}{\milli\ampere}}
          =\qty{0.20}{\ohm},\\
    W &> \frac{\qty{2.5e3}{\ohm\micro\metre}}{\qty{0.20}{\ohm}}
         =\qty{12.5e3}{\micro\metre}=\qty{12.5}{\milli\metre}.
  \end{aligned}
\]
عرض مرزی \qty{12.5}{\milli\metre} در این مدل متناظر با ۲ درصد
افزایش تأخیر است. شرطِ «کمتر از» عرضی بزرگ‌تر از این مرز می‌خواهد.
اگر کلید از چند ترانزیستور موازی ساخته شود، \(W\) مجموع عرض‌های
مؤثر است. برآورد عددی به فرض تغذیه، جریان نماینده و تقریب
افت ولتاژ در مثال کتاب وابسته است.
